\chapter{Casos de uso restantes}
\label{ap:casos-de-uso-restantes}

Este apéndice amplía los siete casos de uso que el capítulo~5 presenta
de forma resumida. Se conservan sus identificadores y su orden para
facilitar la consulta de los diagramas y la matriz de trazabilidad.
Las fichas detallan la interacción de los actores con F.A.M.A. y
distinguen las operaciones de la interfaz de las tareas que requieren
acceso administrativo al nodo.

\section{Gestionar cuenta y perfil}
\label{cu:ses-02}

\noindent\textbf{Identificador.} \texttt{CU\_SES\_02}.
\textbf{Actores.} Investigador y administrador de sistemas.

\noindent\textbf{Objetivo.} Consultar los datos de la cuenta y actualizar
la contraseña de acceso.

\noindent\textbf{Precondiciones.} El usuario tiene una sesión web vigente.
Para cambiar su contraseña debe conocer la contraseña actual.

\subsection*{Flujo principal}

\noindent
\begin{tabular}{@{}p{0.47\linewidth}p{0.48\linewidth}@{}}
    \hline
    \textbf{Actor} & \textbf{Aplicación} \\
    \hline
    1. Abre la vista \emph{Mi cuenta}. &
    2. Muestra el nombre, el correo electrónico y el rol de la sesión. \\
    3. Ingresa la contraseña actual, escribe y confirma la nueva, y solicita el cambio. &
    4. Verifica la contraseña actual y la política de la nueva. Si se cumplen las condiciones, actualiza la credencial e invalida las sesiones de la cuenta. \\
    5. Inicia sesión nuevamente con la contraseña actualizada. &
    6. Verifica las credenciales y establece una nueva sesión. \\
    \hline
\end{tabular}

\subsection*{Flujos alternativos}

\begin{itemize}
    \item \textbf{Solo consulta.} El usuario revisa su identidad y sale
    de la vista sin modificar su contraseña.
    \item \textbf{Cambio rechazado.} Si la contraseña actual es
    incorrecta, la nueva no cumple la política o la confirmación
    difiere, la aplicación informa el error y conserva la credencial.
\end{itemize}

\noindent\textbf{Postcondiciones.} La consulta conserva los datos de la
cuenta. Si se completa el cambio, la nueva contraseña queda vigente y
el acceso continúa mediante una sesión nueva.

\noindent\textbf{Alcance actual.} La identidad se muestra en modo de
consulta. La edición del nombre y la revisión de eventos de seguridad
corresponden a funciones administrativas.

\section{Seleccionar dataset para entrenamiento}
\label{cu:inv-02}

\noindent\textbf{Identificador.} \texttt{CU\_INV\_02}.
\textbf{Actor principal.} Investigador.

\noindent\textbf{Objetivo.} Elegir el conjunto de datos que se utilizará
para configurar un entrenamiento acústico.

\noindent\textbf{Precondiciones.} El investigador tiene una sesión
vigente y acceso a un conjunto de datos registrado como fuente de
entrenamiento.

\subsection*{Flujo principal}

\noindent
\begin{tabular}{@{}p{0.47\linewidth}p{0.48\linewidth}@{}}
    \hline
    \textbf{Actor} & \textbf{Aplicación} \\
    \hline
    1. Abre la vista \emph{Entrenamiento y Métricas}. &
    2. Consulta y muestra los conjuntos de datos visibles para la sesión, con sus cantidades y clases cuando están disponibles. \\
    3. Selecciona el conjunto correspondiente al dominio acústico que desea estudiar. &
    4. Conserva el identificador en el formulario y muestra la preparación que se puede configurar para ese conjunto. \\
    5. Cuando corresponde, elige una versión de las particiones de entrenamiento, validación y prueba. &
    6. Mantiene la versión asociada al conjunto seleccionado para enviarla con la solicitud de entrenamiento. \\
    \hline
\end{tabular}

\subsection*{Flujos alternativos}

\begin{itemize}
    \item \textbf{Catálogo sin opciones.} Si no existen conjuntos
    visibles, el investigador debe disponer de un conjunto registrado
    y autorizado antes de continuar.
    \item \textbf{Fuente local no utilizable.} Si faltan archivos,
    la preparación es incompatible o el acceso cambió, la validación
    de inicio rechaza el trabajo. El investigador puede corregir la
    preparación o elegir otra fuente autorizada.
\end{itemize}

\noindent\textbf{Postcondiciones.} El formulario queda asociado al
conjunto y, cuando se seleccionó, a su versión de particiones. El
inicio y la validación de la fuente pertenecen a \texttt{CU\_INV\_03}.

\section{Consultar historial y métricas}
\label{cu:inv-04}

\noindent\textbf{Identificador.} \texttt{CU\_INV\_04}.
\textbf{Actor principal.} Investigador.

\noindent\textbf{Objetivo.} Revisar los modelos registrados y las
métricas conservadas de su entrenamiento.

\noindent\textbf{Precondiciones.} El investigador tiene una sesión
vigente. Para examinar una ficha debe existir un modelo visible para
la cuenta.

\subsection*{Flujo principal}

\noindent
\begin{tabular}{@{}p{0.47\linewidth}p{0.48\linewidth}@{}}
    \hline
    \textbf{Actor} & \textbf{Aplicación} \\
    \hline
    1. Consulta el historial desde la vista de entrenamiento. &
    2. Recupera los modelos autorizados en PostgreSQL y los presenta desde el más reciente. \\
    3. Abre la ficha del modelo que quiere revisar. &
    4. Presenta los datos registrados, como arquitectura, conjunto de datos, épocas y métricas de acierto (\emph{accuracy}) y pérdida (\emph{loss}). Muestra la ficha del archivo de pesos cuando puede leer sus metadatos. \\
    \hline
\end{tabular}

\subsection*{Flujos alternativos}

\begin{itemize}
    \item \textbf{Historial vacío.} La consulta puede devolver una
    lista sin modelos visibles para la sesión.
    \item \textbf{Ficha incompleta.} Si falta una métrica o el archivo
    de pesos no está disponible para consultar sus metadatos, el
    registro histórico se puede revisar con la información conservada.
\end{itemize}

\noindent\textbf{Postcondiciones.} El investigador obtiene los datos
disponibles para comparar entrenamientos. La consulta mantiene el
estado de activación de los modelos.

\section{Seleccionar modelo activo}
\label{cu:inv-05}

\noindent\textbf{Identificador.} \texttt{CU\_INV\_05}.
\textbf{Actor principal.} Investigador.
\textbf{Actor secundario.} Administrador de sistemas para la designación
del modelo predeterminado.

\noindent\textbf{Objetivo.} Elegir un modelo para analizar una grabación
desde la vista de predicción.

\noindent\textbf{Precondiciones.} La sesión está vigente. El catálogo
de inferencia ofrece un modelo autorizado con disponibilidad de pesos.

\subsection*{Flujo principal}

\noindent
\begin{tabular}{@{}p{0.47\linewidth}p{0.48\linewidth}@{}}
    \hline
    \textbf{Actor} & \textbf{Aplicación} \\
    \hline
    1. Abre el selector de la vista de predicción. &
    2. Consulta el catálogo y presenta los modelos visibles, con sus clases y la disponibilidad observada de sus pesos. \\
    3. Selecciona el modelo correspondiente al tipo de grabación que desea analizar. &
    4. Guarda la elección en la vista y utiliza el identificador seleccionado al solicitar una predicción. \\
    \hline
\end{tabular}

\subsection*{Flujos alternativos}

\begin{itemize}
    \item \textbf{Pesos indisponibles.} El modelo queda fuera de la
    selección cuando el catálogo no confirma sus pesos. Si existe un
    respaldo verificable en GCS, un usuario con acceso puede solicitar
    su restauración y volver a consultar el catálogo.
    \item \textbf{Cambio del predeterminado.} El administrador solicita
    la activación del modelo registrado mediante la API. El servicio
    comprueba los informes de evaluación requeridos y actualiza la
    designación en PostgreSQL. La disponibilidad para inferencia se
    comprueba por separado en el catálogo y en la ejecución.
\end{itemize}

\noindent\textbf{Postcondiciones.} La elección del investigador queda
asociada a su solicitud de predicción (\texttt{CU\_INV\_06}). El
predeterminado del nodo cambia mediante la operación administrativa
de activación.

\section{Gestionar nodos de procesamiento local}
\label{cu:adm-02}

\noindent\textbf{Identificador.} \texttt{CU\_ADM\_02}.
\textbf{Actor principal.} Administrador de sistemas.

\noindent\textbf{Objetivo.} Revisar el estado del nodo y sus recursos
de cómputo antes de operar el flujo acústico.

\noindent\textbf{Precondiciones.} El administrador tiene una sesión
vigente y puede acceder al panel del nodo en funcionamiento.

\subsection*{Flujo principal}

\noindent
\begin{tabular}{@{}p{0.47\linewidth}p{0.48\linewidth}@{}}
    \hline
    \textbf{Actor} & \textbf{Aplicación} \\
    \hline
    1. Abre el panel general y consulta el estado del nodo. &
    2. Recopila la disponibilidad de CUDA, los datos de CPU y RAM del equipo y, si detecta GPU, la memoria gráfica observada por el proceso. \\
    3. Revisa el dispositivo detectado, la memoria y el nivel de carga para decidir si inicia un trabajo. &
    4. Muestra la telemetría disponible y actualiza la lectura mediante consultas periódicas. \\
    \hline
\end{tabular}

\subsection*{Flujos alternativos}

\begin{itemize}
    \item \textbf{CUDA no detectado.} La telemetría identifica el
    entorno de CPU. El administrador puede revisar el acceso de la
    aplicación a la GPU en el equipo anfitrión.
    \item \textbf{Telemetría indisponible.} Si el panel no obtiene una
    lectura, el administrador comprueba la comunicación con el
    backend y vuelve a consultar.
\end{itemize}

\noindent\textbf{Postcondiciones.} El administrador dispone de un
diagnóstico del nodo en el momento de la consulta. Los cambios de
controladores o de configuración de hardware se realizan en el
entorno de despliegue. El panel se utiliza para comprobar su efecto.

\section{Gestionar recursos cloud}
\label{cu:adm-03}

\noindent\textbf{Identificador.} \texttt{CU\_ADM\_03}.
\textbf{Actor principal.} Administrador de sistemas.
\textbf{Actor secundario.} Google Cloud Storage (GCS).

\noindent\textbf{Objetivo.} Comprobar la conexión del nodo con el
\emph{bucket} configurado y consultar su inventario general.

\noindent\textbf{Precondiciones.} El administrador tiene una sesión
vigente. La consulta requiere un \emph{bucket} configurado, credenciales
del servicio y conectividad hacia GCS.

\subsection*{Flujo principal}

\noindent
\begin{tabular}{@{}p{0.47\linewidth}p{0.48\linewidth}@{}}
    \hline
    \textbf{Actor} & \textbf{Aplicación} \\
    \hline
    1. Abre la vista de ingesta y consulta el estado del almacenamiento remoto. &
    2. Comprueba el rol de la cuenta y solicita a GCS el listado de objetos del \emph{bucket} configurado. \\
    3. Revisa el estado de conexión y el tamaño del inventario. &
    4. Informa el \emph{bucket}, el total de objetos y su tamaño acumulado cuando la consulta se completa. \\
    \hline
\end{tabular}

\subsection*{Flujos alternativos}

\begin{itemize}
    \item \textbf{Conexión fallida.} Ante un problema de red, identidad,
    permisos o configuración del \emph{bucket}, el servicio informa
    que la conexión está indisponible.
    \item \textbf{Configuración inicial o corrección.} Durante la
    instalación, el operador autorizado puede definir el proyecto y
    el \emph{bucket} en el asistente y ejecutar sus comprobaciones.
    Los cambios de credenciales de GCS se aplican en el entorno del
    nodo. Tras un ajuste, se vuelve a comprobar la conexión.
\end{itemize}

\noindent\textbf{Postcondiciones.} Queda disponible el resultado de la
consulta de conexión y, si tuvo éxito, el inventario observado.
Esta operación aporta evidencia para \texttt{RF\_01}.

\section{Controlar ejecución de procesos}
\label{cu:adm-05}

\noindent\textbf{Identificador.} \texttt{CU\_ADM\_05}.
\textbf{Actor principal.} Administrador de sistemas.

\noindent\textbf{Objetivo.} Solicitar la detención de un entrenamiento
en curso y comprobar su estado final.

\noindent\textbf{Precondiciones.} El administrador tiene una sesión
vigente y el nodo registra un entrenamiento activo.

\subsection*{Flujo principal}

\noindent
\begin{tabular}{@{}p{0.47\linewidth}p{0.48\linewidth}@{}}
    \hline
    \textbf{Actor} & \textbf{Aplicación} \\
    \hline
    1. Revisa el progreso del entrenamiento y solicita detenerlo. &
    2. Comprueba el permiso y el estado del trabajo. Registra una señal de detención cooperativa e informa que la solicitud está en curso. \\
    3. Continúa consultando el estado de la ejecución. &
    4. El proceso atiende la señal en un punto de control y el servicio refleja la detención cuando esta concluye. \\
    \hline
\end{tabular}

\subsection*{Flujos alternativos}

\begin{itemize}
    \item \textbf{Trabajo ya finalizado.} Si ya no existe un
    entrenamiento activo, el servicio devuelve el estado vigente e
    informa que no hay un trabajo que detener.
    \item \textbf{Detención pendiente.} La respuesta \texttt{stopping}
    confirma la recepción de la solicitud. La detención se confirma
    cuando la consulta de progreso devuelve \texttt{stopped}.
\end{itemize}

\noindent\textbf{Postcondiciones.} Cuando se confirma la detención, el
trabajo deja de ejecutarse. Los modelos ya registrados se conservan.
El modelo que estaba en entrenamiento puede quedar sin publicar.

\noindent\textbf{Alcance actual.} La detención de entrenamiento es
cooperativa. La cancelación de trabajos de importación ZIP se gestiona
por separado desde su módulo, con los permisos del titular de la
importación.

% Fichas ampliadas de las tablas 5.2, 5.4, 5.6, 5.7, 5.11, 5.12 y 5.14.
% Contraste con backend/app/main.py, backend/app/installation.py,
% backend/app/routes/auth.py, backend/app/routes/checkpoints.py,
% backend/app/services/training.py y backend/app/services/ingestion.py.
% Vistas de referencia: AccountView, TrainingView, PredictionView,
% DashboardView e IngestionView en frontend/components/.
